\documentclass[11pt]{article}

%------------------------- Packages -------------------------
\usepackage[margin=1in]{geometry}
\usepackage{amsmath, amssymb}
\usepackage{graphicx}
\usepackage{listings}
\usepackage[hidelinks]{hyperref}
\usepackage{enumitem}
\usepackage{indentfirst}
\usepackage{titling}

% Figures are resolved from results/figures/ (this file lives in results/reports/).
\graphicspath{{../figures/}}

%------------------------- Paragraph indentation -------------------------
\setlength{\parindent}{2em}
\setlength{\parskip}{0pt plus 1pt}
\setlength{\droptitle}{-1.75cm}

%------------------------- Code style -------------------------
\lstset{
    basicstyle=\ttfamily\small,
    breaklines=true,
    frame=single
}

\title{\textbf{Lab1 Report}}
\author{}
\date{}

\begin{document}
\maketitle

\section*{Task 1 \& Task 2}
After setting the \texttt{GlobalFrequency = 2GHz} and adding the new metric, the result \texttt{network\_stats.txt} contains:

\begin{lstlisting}[caption={Example network\_stats.txt contents}]
packets_injected = 3165                       (Count)
packets_received = 3162                       (Count)
average_packet_queueing_latency = 2           (Tick/Count)
average_packet_network_latency = 13.557559    (Tick/Count)
average_packet_latency = 15.557559            (Tick/Count)
average_hops = 5.269450                       (Count/Count)
reception_rate = 0.004941                     (packets/node/cycle)
\end{lstlisting}
\begin{figure}[h]
    \centering
    \includegraphics[width=0.9\textwidth]{Snipaste_2026-08-23_21-40-37.png}
    \caption{Screenshot of \texttt{network\_stats.txt}}
\end{figure}
Note:

The Reception Rate of $0.004941$ packets/node/cycle is only about half of the configured injection rate of $0.01$ because the simulation effectively executed roughly $5000$ cycles instead of the $10000$ cycles implied by \texttt{--sim-cycles}. The root cause is a unit mismatch in the termination condition of the synthetic traffic generator. In gem5, \texttt{GlobalFrequency} defines the duration of one tick; setting it to 2GHz makes one tick equal to $0.5$ ns. Meanwhile, the default \texttt{--sys-clock} is 1GHz, so one CPU cycle lasts $1$ ns, i.e., exactly $2$ ticks. However, the generator terminates the simulation inside \texttt{GarnetSyntheticTraffic::tick()} with the check \texttt{curTick() >= simCycles}, which compares the current absolute time measured in ticks against the \texttt{--sim-cycles} value that is nominally expressed in cycles. Consequently, when \texttt{--sim-cycles = 10000} is given, the simulation stops after only $10000$ ticks, which corresponds to $10000 / 2 = 5000$ actual clock cycles rather than $10000$.

\section*{Task 3}

\subsection*{Q1}
What input parameter options are available for the ``garnet\_synth\_traffic.py''? Briefly describe their usage. Where (in which files) are the default parameters defined? (Tips: use -h option)\\

\begin{lstlisting}[caption={Options for garnet\_synth\_traffic.py}]
  -h, --help            show this help message and exit
  -n NUM_CPUS, --num-cpus NUM_CPUS
  --sys-voltage SYS_VOLTAGE
                        Top-level voltage for blocks running at system power supply
  --sys-clock SYS_CLOCK
                        Top-level clock for blocks running at system speed
  --list-mem-types      List available memory types
  --mem-type {CfiMemory,DDR3_1600_8x8,DDR3_2133_8x8,DDR4_2400_16x4,DDR4_2400_4x16,DDR4_2400_8x8,DDR5_4400_4x8,DDR5_6400_4x8,DDR5_8400_4x8,DRAMInterface,GDDR5_4000_2x32,HBM_1000_4H_1x128,HBM_1000_4H_1x64,HBM_2000_4H_1x64,HMC_2500_1x32,LPDDR2_S4_1066_1x32,LPDDR3_1600_1x32,LPDDR5_5500_1x16_8B_BL32,LPDDR5_5500_1x16_BG_BL16,LPDDR5_5500_1x16_BG_BL32,LPDDR5_6400_1x16_8B_BL32,LPDDR5_6400_1x16_BG_BL16,LPDDR5_6400_1x16_BG_BL32,NVMInterface,NVM_2400_1x64,QoSMemSinkInterface,SimpleMemory,WideIO_200_1x128}
                        type of memory to use
  --mem-channels MEM_CHANNELS
                        number of memory channels
  --mem-ranks MEM_RANKS
                        number of memory ranks per channel
  --mem-size MEM_SIZE   Specify the physical memory size (single memory)
  --enable-dram-powerdown
                        Enable low-power states in DRAMInterface
  --mem-channels-intlv MEM_CHANNELS_INTLV
                        Memory channels interleave
  --memchecker
  --external-memory-system EXTERNAL_MEMORY_SYSTEM
                        use external ports of this port_type for caches
  --tlm-memory TLM_MEMORY
                        use external port for SystemC TLM cosimulation
  --caches
  --l2cache
  --num-dirs NUM_DIRS
  --num-l2caches NUM_L2CACHES
  --num-l3caches NUM_L3CACHES
  --l1d_size L1D_SIZE
  --l1i_size L1I_SIZE
  --l2_size L2_SIZE
  --l3_size L3_SIZE
  --l1d_assoc L1D_ASSOC
  --l1i_assoc L1I_ASSOC
  --l2_assoc L2_ASSOC
  --l3_assoc L3_ASSOC
  --cacheline_size CACHELINE_SIZE
  --ruby
  -m TICKS, --abs-max-tick TICKS
                        Run to absolute simulated tick specified including ticks from a restored checkpoint
  --rel-max-tick TICKS  Simulate for specified number of ticks relative to the simulation start tick (e.g. if
                        restoring a checkpoint)
  --maxtime MAXTIME     Run to the specified absolute simulated time in seconds
  -P PARAM, --param PARAM
                        Set a SimObject parameter relative to the root node. An extended Python multi range slicing
                        syntax can be used for arrays. For example: 'system.cpu[0,1,3:8:2].max_insts_all_threads =
                        42' sets max_insts_all_threads for cpus 0, 1, 3, 5 and 7 Direct parameters of the root
                        object are not accessible, only parameters of its children.
  --synthetic {uniform_random,tornado,bit_complement,bit_reverse,bit_rotation,neighbor,shuffle,transpose}
  -i I, --injectionrate I
                        Injection rate in packets per cycle per node. Takes decimal value between 0 to 1 (eg.
                        0.225). Number of digits after 0 depends upon --precision.
  --precision PRECISION
                        Number of digits of precision after decimal point for injection rate
  --sim-cycles SIM_CYCLES
                        Number of simulation cycles
  --num-packets-max NUM_PACKETS_MAX
                        Stop injecting after --num-packets-max. Set to -1 to disable.
  --single-sender-id SINGLE_SENDER_ID
                        Only inject from this sender. Set to -1 to disable.
  --single-dest-id SINGLE_DEST_ID
                        Only send to this destination. Set to -1 to disable.
  --inj-vnet {-1,0,1,2}
                        Only inject in this vnet (0, 1 or 2). 0 and 1 are 1-flit, 2 is 5-flit. Set to -1 to inject
                        randomly in all vnets.
  --ruby-clock RUBY_CLOCK
                        Clock for blocks running at Ruby system's speed
  --access-backing-store
                        Should ruby maintain a second copy of memory
  --ports PORTS         used of transitions per cycle which is a proxy for the number of ports.
  --numa-high-bit NUMA_HIGH_BIT
                        high order address bit to use for numa mapping. 0 = highest bit, not specified = lowest bit
  --interleaving-bits INTERLEAVING_BITS
                        number of bits to specify interleaving in directory, memory controllers and caches. 0 = not
                        specified
  --xor-low-bit XOR_LOW_BIT
                        hashing bit for channel selectionsee MemConfig for explanation of the defaultparameter. If
                        set to 0, xor_high_bit is alsoset to 0.
  --recycle-latency RECYCLE_LATENCY
                        Recycle latency for ruby controller input buffers
  --topology TOPOLOGY   check configs/topologies for complete set
  --mesh-rows MESH_ROWS
                        the number of rows in the mesh topology
  --network {simple,garnet}
                        'simple'|'garnet' (garnet2.0 will be deprecated.)
  --router-latency ROUTER_LATENCY
                        number of pipeline stages in the garnet router. Has to be >= 1. Can be over-ridden on a per
                        router basis in the topology file.
  --link-latency LINK_LATENCY
                        latency of each link the simple/garnet networks. Has to be >= 1. Can be over-ridden on a
                        per link basis in the topology file.
  --link-width-bits LINK_WIDTH_BITS
                        width in bits for all links inside garnet.
  --vcs-per-vnet VCS_PER_VNET
                        number of virtual channels per virtual network inside garnet network.
  --routing-algorithm ROUTING_ALGORITHM
                        routing algorithm in network. 0: weight-based table 1: XY (for Mesh. see
                        garnet/RoutingUnit.cc) 2: Custom (see garnet/RoutingUnit.cc
  --network-fault-model
                        enable network fault model: see src/mem/ruby/network/fault_model/
  --garnet-deadlock-threshold GARNET_DEADLOCK_THRESHOLD
                        network-level deadlock threshold.
  --simple-physical-channels
                        SimpleNetwork links uses a separate physical channel for each virtual network
\end{lstlisting}

The traffic options are defined in \texttt{configs/example/garnet\_synth\_traffic.py}: \texttt{synthetic}, \texttt{injectionrate}, \texttt{precision}, \texttt{sim-cycles}, \texttt{num-packets-max}, \texttt{single-sender-id}, \texttt{single-dest-id} and \texttt{inj-vnet}. 

The other options are defined in \texttt{configs/common/Options.py}.

\subsection*{Q2}
What is the unit of "sim-cycles"? What is the unit of "router-latency" and "link-latency". What is the relationship between Tick and Cycle?\\

\texttt{sim-cycles}: cycles. 

\texttt{router-latency}: cycles (number of router pipeline stages). 

\texttt{link-latency}: cycles (flit traversal latency on a link). 

A Tick is the absolute time unit of the simulator, fixed by \texttt{GlobalFrequency} (default $1\ \text{tick}=1\ \text{ps}$); a Cycle is the period of a particular clock domain. Their relationship is
\[
1\ \text{cycle (ticks)} = \frac{f_{\text{tick}}}{f_{\text{clock}}},
\]
where $f_{\text{tick}}$ is the global tick frequency and $f_{\text{clock}}$ is the clock-domain frequency. For example, with $1\ \text{tick}=1\ \text{ps}$ and a $1\ \text{GHz}$ domain, $1\ \text{cycle}=1000\ \text{ticks}$.

\subsection*{Q3}
What is the unit of "injectionrate"?\\

Packets per cycle per node (packets/cycle/node) or (count/cycle/count).

\subsection*{Q4}
Where are the GarnetNetworkInterface and GarnetRouter defined?\\

\texttt{GarnetNetworkInterface} in \texttt{src/mem/ruby/network/garnet/NetworkInterface.hh/.cc}.

\texttt{GarnetRouter} in \texttt{src/mem/ruby/network/garnet/Router.hh/.cc}.


\subsection*{Q5}
In which module(s), packets are generated and injected to the network? In which module(s), packets are buffered during transmission? In which module(s), the program determines if packets can be sent downstream.\\

Generation and injection: \texttt{GarnetSyntheticTraffic} and \texttt{NetworkInterface}.

Buffering: \texttt{NetworkInterface} (\texttt{niOutVcs}), \texttt{Router} (\texttt{InputUnit}/\texttt{OutputUnit}), and \texttt{NetworkLink} (\texttt{linkBuffer}).

Downstream decision: \texttt{Router::wakeup()} (\texttt{SwitchAllocator}/\texttt{OutputUnit}).

\end{document}
